\section{Interstate 2.0: Managed Lanes and Incident Analysis}
\label{sec:i2scenario}

\subsection{Incident Simulation: Donner Pass Closure}

A Donner Pass multi-day closure is simulated by removing the I-80 Sierra Nevada edge from the directed interstate graph (Section~\ref{sec:methods}) and recomputing the shortest path from the NYC source node to the LA sink node. The forced detour routes all NY--LA traffic via I-80 to Sacramento, south on CA-99 or CA-5 to the I-5/CA-58 junction near Bakersfield, and then to I-40 at Barstow for the Los Angeles approach. This is not the textbook I-40 alternate (which would require diverting east to Reno and entering I-40 from there); it is the emergency routing that Caltrans activates for Donner closures, using the Central Valley as a north-south bypass.

The practical alternate adds 310 miles to the trip distance relative to the Donner route. At 55 mph:
\begin{equation}
  \Delta T_{\text{detour}} = \frac{310 \text{ mi}}{55 \text{ mph}} = 5.6 \text{ hr}
\end{equation}

The daily economic cost of a Donner closure, applied to the 8,000 trucks per day that transit the Sierra Nevada on I-80 \citep{Caltrans2023}:
\begin{equation}
  C_{\text{day}} = 8{,}000 \times 5.6 \times \$225 = \$10.1\text{M/day} \approx \$9.9\text{M/day}
\end{equation}

Over the 50 annual closure events with a mean duration of 18 hours:
\begin{equation}
  C_{\text{detour,annual}} = \$9.9\text{M/day} \times \frac{18}{24} \times 50 = \$371\text{M}
\end{equation}

This detour cost does not capture the trucks that are in transit at Donner when the closure begins and must wait for reopening rather than diverting. The 18-hour mean wait cost for detained trucks:
\begin{equation}
  C_{\text{wait,annual}} = 8{,}000 \times 0.25 \times 18 \times \$225 \times 50 = \$405\text{M}
\end{equation}
where 0.25 is the estimated fraction of daily truck volume that is in the immediate Donner vicinity (within 50 miles in either direction) when a closure activates, based on Caltrans closure notification data \citep{Caltrans2023}. Combined Donner closure cost: approximately \$776 million per year in direct operating cost, with additional shipper-side delay costs (missed delivery windows, penalty payments, inventory carrying costs) that bring the total toward \$1 billion annually.

Max-flow analysis confirms that Donner closure reduces NY--LA corridor max-flow by 91,200 vpd to zero on the primary path; the alternate path has a binding constraint at the I-5/CA-58 junction at approximately 73,000 vpd. The network does not fail completely, but the primary path's max-flow is eliminated.

\subsection{Incident Simulation: Dallas Interchange Failure}

A major incident at the Dallas I-45/I-35 interchange (the most likely single point of failure for the HOU--CHI corridor, based on FHWA incident frequency data) forces HOU--CHI traffic onto an alternate: south on I-45 to Huntsville, west on US-190 to I-35, bypassing the Dallas core. This bypass adds approximately 85 miles and 1.5 hours. For incidents severe enough to close I-35 through Dallas entirely (a rare but documented event for multi-vehicle accidents in managed lanes), the alternate routes traffic via I-30 east to Memphis and then I-55 north---adding approximately 200 miles.

Mean time to clearance for major incidents on T1 corridors: 4.2 hours \citep{FHWA_freight2023}. At 3,200 trucks per day affected:
\begin{equation}
  C_{\text{Dallas incident}} = 3{,}200 \times 4.2 \times \$225 \times (\text{annual frequency}) = \$3.0\text{M per event}
\end{equation}
ATRI records approximately 180 major delay events per year at the Dallas I-45/I-35 node \citep{ATRI_costs2024}, yielding approximately \$540 million per year in incident-induced delay cost at Dallas alone.

\subsection{Interstate 2.0 Managed Freight Lane Specification}

The Interstate 2.0 (I2.0) managed freight lane design applied to the NY--LA and HOU--CHI corridors in this scenario has the following specification:

\begin{itemize}
  \item \textbf{Lane count}: 2 dedicated freight-only lanes per direction, access-controlled, no passenger vehicles
  \item \textbf{Separation}: Concrete barrier separation from general-purpose lanes; dedicated on/off ramps at freight interchange nodes
  \item \textbf{Design speed}: 75 mph (sustained operating speed 65 mph, accounting for HOS deceleration and acceleration events)
  \item \textbf{Grade standards}: Maximum 3 percent sustained grade on new alignments; existing Donner Pass grade (4\%) excepted but mitigated by truck climbing lanes
  \item \textbf{Weather hardening}: Variable message signs, real-time chain control status, and preemptive closure protocols reducing mean Donner closure from 18 hours to an estimated 12 hours (via earlier proactive closure and faster reopening with dedicated equipment)
  \item \textbf{Design V/C}: 0.70 (LOS B; uncongested flow with reserve capacity for demand surge)
\end{itemize}

\subsection{NY--LA: I2.0 Transit and PTI}

Under I2.0 managed lanes at 65 mph sustained:
\begin{equation}
  T_{\text{NY-LA, I2.0}} = \frac{2{,}800 \text{ mi}}{65 \text{ mph}} \div 11 \text{ hr/day} = 43.1 \text{ hr} \div 11 = 3.9 \text{ days}
\end{equation}
Reduced stops from coordinated rest facilities at I2.0 freight nodes (reducing average stop overhead from 45 minutes to 20 minutes per rest event, with 3 rest events on a transcontinental trip) saves an additional 1.25 hours, yielding an effective transit of approximately 3.5 days---consistent with the abstract figure.

The I2.0 PTI target of 1.15 applies to the managed lanes. Applying this to the free-flow transit time:
\begin{align}
  T_{\text{ff, I2.0}} &= \frac{2{,}800}{65} = 43.1 \text{ hr} \\
  T_{95, \text{I2.0}} &= 43.1 \times 1.15 = 49.6 \text{ hr}
\end{align}
A shipper committing to 95-percent on-time delivery within a 48-hour window on a 43-hour free-flow trip requires PTI $\leq 48/43 = 1.116$. The I2.0 target of PTI 1.15 achieves this with a 2-hour margin (50 hours 95th-percentile vs.\ 48-hour commitment). The 48-hour SLA is feasible.

Compare to current performance: PTI 1.86 on a 43-hour free-flow trip requires a commitment window of $43 \times 1.86 = 80$ hours for 95-percent confidence. The I2.0 improvement from 80-hour to 48-hour commitment represents a 40-percent reduction in the required SLA buffer---a commercially transformative change for shippers managing just-in-time inventory.

\subsection{HOU--CHI with I-69: I2.0 Transit and PTI}

With I-69 completion (870-mile direct corridor) and I2.0 managed lanes at 65 mph:
\begin{equation}
  T_{\text{HOU-CHI, I2.0}} = \frac{870 \text{ mi}}{65 \text{ mph}} \div 11 \text{ hr/day} = 13.4 \text{ hr} \div 11 = 1.22 \text{ days} \approx 1.5 \text{ days}
\end{equation}
The 13.4-hour free-flow transit is achievable in a single HOS driving day for a team driver, or in two driving windows for a solo driver. At PTI 1.15, the 95th-percentile trip takes $13.4 \times 1.15 = 15.4$ hours. A shipper can offer a 24-hour SLA from Houston to Chicago with 95-percent confidence---compared to the current 48-hour quote under the three-hop route. Same-day delivery from Houston to Chicago becomes feasible for freight departing before 8~am.

\subsection{Combined Reliability Cost and I2.0 Benefit}

\begin{table}[ht]
\centering
\caption{Before/after comparison: I2.0 scenario applied to NY--LA and HOU--CHI corridors}
\label{tab:i2_comparison}
\begin{tabular}{lllll}
\toprule
Metric & NY--LA Current & NY--LA I2.0 & HOU--CHI Current & HOU--CHI I2.0+I-69 \\
\midrule
Distance (mi)          & 2,800 & 2,800 & 1,160 & 870 \\
Transit time (days)    & 4.5   & 3.5   & 1.92  & 1.22 \\
Corridor PTI           & 1.86  & 1.15  & $\sim$1.45 & 1.15 \\
SLA window (95\%)      & 80 hr & 48 hr & $\sim$54 hr & 24 hr \\
Annual reliability cost & \$5.7B & --- & \$2.5B & --- \\
\bottomrule
\end{tabular}
\end{table}

The \$5.7 billion annual reliability cost for NY--LA is computed from:
\begin{itemize}
  \item Donner closure events: \$776 million (direct operating cost)
  \item Bay Area congestion delay (V/C 1.86 on 8,000 trucks/day, 260 peak days/year, mean 45-minute delay): \$2.16 billion
  \item Shipper-side costs (inventory buffer carrying cost at 20\% annual rate on \$12B of goods in the reliability buffer zone): \$2.4 billion
  \item Remaining segment delay and uncertainty costs: \$0.36 billion
\end{itemize}
The \$2.5 billion HOU--CHI figure is derived in Section~\ref{sec:houtochi}. Combined: \$8.2 billion per year in addressable freight reliability cost \citep{ATRI_costs2024}.

I2.0 investment does not eliminate all of this cost. Managed lanes reduce congestion delay and improve PTI; they do not eliminate weather-related Donner closures, though hardened infrastructure and proactive management reduce both frequency and duration. The addressable fraction---the portion attributable to congestion, distance penalty, and interchange variance---is approximately 70 percent of the combined total, or \$5.7 billion. The remaining 30 percent requires physical infrastructure (Donner tunnel or bore, I-69 gap completion) that is beyond managed-lane retrofit.
